\paragraph{Cluster A Step 15: minimal anomaly-support test.}

The gauge group \(SU(3)\times SU(2)\times U(1)\) is held fixed.  The finite
enumeration contains \(66\) multiplet types:
\[
SU(3)\in\{1,3,\bar 3\},\qquad SU(2)\in\{1,2\},\qquad
Y\in \frac{1}{6}\{-6,-4,-3,-2,-1,0,1,2,3,4,6\},
\]
with at most \(5\) distinct Weyl multiplet types per support.

The real anomaly filter finds \(1161\) raw anomaly-free supports.
After quotienting vector-like pairs, overall hypercharge normalization, and
trivial neutral singlets, \(28\) irreducible chiral supports remain.

The declared minimality measure is:
\[
\text{score}(S)=(\#\text{ irreducible chiral multiplets},\ \text{total Weyl components}).
\]
The SM one-generation support is present and anomaly-free, with score
\((5,15)\), but the minimal score in the enumeration is
\((4,12)\). Thus the SM support is not selected
by anomaly freedom plus this minimality measure.

\paragraph{Verdict.}
Minimality alone is insufficient in this finite enumeration.  This is a typed
no-go for the scoped selection-law candidate, not a wall.  The next principle to
test is an electroweak quark-lepton participation or observed-charge sector
condition.

\paragraph{Grade.}
Finite-carrier diagnostic, conditional on the fixed gauge group and bounded
enumeration.  No physical value, mechanism, or frame-transfer closure is supplied.
